\documentclass{ieeeaccess}

\def\BibTeX{{\rm B\kern-.05em{\sc i\kern-.025em b}\kern-.08em
    T\kern-.1667em\lower.7ex\hbox{E}\kern-.125emX}}

\begin{document}
\history{Date of publication xxxx 00, 0000, date of current version xxxx 00, 0000.}
\doi{10.1109/ACCESS.2026.DOI}

\title{An LLM-Driven MCP Framework for Automated DoDAF-Compliant Architecture Model Generation: A Case Study on Smart City Emergency Systems}

\author{\IEEEauthorblockN{Weigang Liu\IEEEauthorrefmark{1}, Xing Zhou\IEEEauthorrefmark{1}, 
Zhongjie Qi\IEEEauthorrefmark{1}, Zhenfeng Yan\IEEEauthorrefmark{1}, 
Xuepan Gao\IEEEauthorrefmark{1}\IEEEauthorrefmark{2}, Chuan Geng\IEEEauthorrefmark{1}, and Haiyang Ren\IEEEauthorrefmark{1}}

\IEEEauthorblockA{\IEEEauthorrefmark{1}The Network and Communication Research Institute of 
China Electronics Technology Group Corporation, Shijiazhuang, Hebei 050081, China}

\IEEEauthorblockA{\IEEEauthorrefmark{2}Corresponding author: Xuepan Gao (e-mail: TBD)}
}

\tfootnote{This work was supported by the Network and Communication Research Institute of China Electronics Technology Group Corporation.}

\markboth
{Liu \headeretal: LLM-Driven MCP Framework for Automated DoDAF Architecture Model Generation}
{Liu \headeretal: LLM-Driven MCP Framework for Automated DoDAF Architecture Model Generation}

\begin{abstract}
Department of Defense Architecture Framework (DoDAF) provides a standardized approach for military system-of-systems architecture modeling, but manual construction of DoDAF-compliant models remains labor-intensive, error-prone, and requires deep domain expertise. Prior automation approaches---including Model-Driven Architecture (MDA) transformations and ontology-based reasoning---are inherently rule-driven and require formally structured inputs, rendering them ineffective when faced with ambiguous, evolving, or informally stated early-stage requirements. This paper proposes an end-to-end framework that leverages Large Language Models (LLMs) as intelligent agents to autonomously orchestrate Model Context Protocol (MCP)-enabled Enterprise Architecture (EA) tools for generating DoDAF-compliant architecture models. Critically, MCP bridges the \emph{action space gap} between LLMs and professional EA tools (e.g., Sparx Systems Enterprise Architect), enabling LLMs to directly manipulate model repositories rather than merely producing textual diagram descriptions such as PlantUML. The framework employs a multi-agent workflow comprising a Requirement Agent, a Planner Agent, an Executer Agent, and a Verifier Agent, which collectively interpret natural language requirements, autonomously select relevant DoDAF viewpoints through a formalized keyword-weight mapping mechanism, and invoke MCP tools to create Systems Modeling Language (SysML) diagrams, model elements, and connectors. We evaluate the framework through three rounds of experiments using a Smart City Emergency Command System case study with DeepSeek-V4 as the LLM backend. Results demonstrate an average end-to-end automated pipeline time of 4~min~50~s per full architecture---including generation, Verifier Agent detection, and automated correction---with 61.3 MCP tool calls, 100\% tool-call reliability (zero failures across 184 invocations), and an average of 9.0 SysML diagrams generated per session. Crucially, we distinguish between pipeline reliability and semantic correctness: the Verifier Agent identified and corrected an average of 3.7 consistency issues per round (81\% reduction), and post-verification models achieved a domain-expert-assessed correctness score of 92.3\%. This work contributes the first LLM-MCP-DoDAF integrated pipeline with a formalized viewpoint selection mechanism and closed-loop quality verification, and provides empirical evidence separating efficiency, reliability, and correctness---three dimensions that prior work has often conflated.
\end{abstract}

\begin{IEEEkeywords}
DoDAF, Large Language Models, Model Context Protocol, Enterprise Architecture, SysML, Automated Model Generation, Smart City Emergency Systems
\end{IEEEkeywords}

\maketitle

\section{Introduction}
\label{sec:introduction}

The Department of Defense Architecture Framework (DoDAF) serves as the cornerstone for military system-of-systems architecture modeling, providing a comprehensive set of 52 model products across eight viewpoints~\cite{DoDAF2010}. However, the construction of DoDAF-compliant architecture models remains a significant bottleneck in defense system engineering. Domain experts must manually translate operational requirements into formal architecture representations using specialized Enterprise Architecture (EA) tools such as Sparx Systems Enterprise Architect (EA), which demands deep familiarity with both DoDAF semantics and tool-specific operations. This manual process is time-consuming, inconsistent across practitioners, and prone to omissions in viewpoint coverage.

\subsection{Why Prior Automation Approaches Fall Short}

Existing approaches to DoDAF automation---including Model-Driven Architecture (MDA) transformations~\cite{mda_dodaf}, ontology-based reasoning~\cite{ontology_dodaf}, and template-based generation~\cite{template_dodaf}---share a fundamental limitation: they are \emph{hard-coded or rule-driven} systems that require formally structured, unambiguous input specifications. In practice, however, defense system requirements in early concept development phases are inherently ambiguous, incomplete, and continuously evolving. These approaches exhibit near-zero tolerance for linguistic vagueness; a single imprecise term or omitted relationship can cause the entire transformation pipeline to fail or, worse, produce syntactically valid but semantically meaningless output. Furthermore, these methods cannot perform \emph{semantic bridging}---the ability to map an informally stated concept (e.g., ``the emergency command center coordinates all responding agencies'') to a formally defined architectural construct (e.g., a SysML Block with specific ports and interfaces).

\subsection{Why LLMs and MCP Are Necessary}

Recent advances in Large Language Models (LLMs) such as GPT-4 and DeepSeek-V4~\cite{deepseek} have demonstrated remarkable capabilities in natural language understanding, contextual reasoning, and tool orchestration~\cite{llm_survey}. LLMs address the semantic gap that rule-driven systems cannot: they can interpret ambiguous requirements, resolve contextual ambiguities, and perform the ``fuzzy-to-precise'' mapping from natural language to formal architectural constructs. This represents a qualitative leap from brittle, rule-based parsing to robust, context-aware semantic interpretation.

However, LLMs alone are insufficient for architecture model generation. Prior LLM-based modeling approaches~\cite{gpt4_uml,llm_requirements} can only produce textual diagram descriptions (e.g., PlantUML code or Mermaid syntax), which users must then manually import into professional EA tools. This \emph{action space gap}---the inability of LLMs to directly manipulate the underlying model repositories of tools like Enterprise Architect---represents a critical bottleneck. LLMs are confined to the action space of text generation; they cannot create, modify, or connect elements within a professional EA tool's database. The Model Context Protocol (MCP), introduced by Anthropic as an open standard for AI-tool interoperability~\cite{mcp_spec}, fundamentally resolves this gap. MCP provides a standardized interface through which LLMs can invoke tool operations (create packages, instantiate elements, draw diagrams, establish connectors) directly on the EA model repository. In essence, MCP expands the LLM's action space from text generation to tool-mediated model construction.

\subsection{Contributions}

This paper addresses the research question: \textit{Can LLMs, guided by MCP, automatically generate DoDAF-compliant architecture models from natural language requirements with acceptable quality and efficiency?} Our contributions are fourfold:

\begin{enumerate}
    \item \textbf{End-to-End Generation Framework}: We present the first integrated pipeline that chains natural language requirements through multi-agent LLM reasoning to MCP-enabled EA tool operations for automated DoDAF architecture generation, selecting key viewpoints rather than exhaustively covering all 52 model products.
    \item \textbf{Formalized Viewpoint Selection}: We propose a context-aware viewpoint selection mechanism based on semantic keyword-weight mapping (Algorithm~1), making the LLM's architectural decision-making transparent, configurable, and reproducible.
    \item \textbf{MCP-Enabled Action Space}: We demonstrate that MCP bridges the action space gap between LLMs and professional EA tools, enabling direct repository manipulation rather than textual code generation. We provide concrete examples of MCP tool call/response JSON exchanges.
    \item \textbf{Empirical Evaluation with Quality Assessment}: We conduct three rounds of experiments using a Smart City Emergency Command System case study with DeepSeek-V4, evaluating not only efficiency (4~min~50~s end-to-end, including Verifier detection and auto-correction) and reliability (100\% tool-call success) but also semantic correctness (92.3\% expert score) through a Verifier Agent and domain expert assessment. An ablation study quantifies the Verifier Agent's contribution (81\% reduction in consistency issues).
\end{enumerate}

The remainder of this paper is organized as follows. Section~\ref{sec:related} reviews related work in NLP-assisted modeling, MCP protocol design, and DoDAF automation, highlighting the limitations of prior approaches. Section~\ref{sec:methodology} presents our framework architecture, MCP tool design with concrete examples, multi-agent workflow, and formalized viewpoint selection. Section~\ref{sec:casestudy} details the Smart City Emergency Command System case study with illustrative generated diagrams. Section~\ref{sec:evaluation} provides experimental results across efficiency, reliability, and semantic correctness dimensions, including an ablation study. Section~\ref{sec:conclusion} concludes with limitations and future directions.

\section{Related Work}
\label{sec:related}

\subsection{NLP-Assisted Architecture Modeling}

The application of Natural Language Processing (NLP) to software and system modeling has gained traction in recent years. Early efforts focused on extracting UML class diagrams from textual requirements using rule-based parsers and template matching~\cite{uml_nlp}. With the advent of deep learning, sequence-to-sequence models were applied to generate structured model elements from requirement descriptions~\cite{deep_modeling}. However, these approaches typically target a single diagram type (e.g., class diagrams or sequence diagrams) and lack the multi-viewpoint coordination required by enterprise architecture frameworks like DoDAF. More fundamentally, they inherit the same brittleness as MDA approaches: their rule-based or learned mappings fail when inputs deviate from expected patterns.

Recent work has explored using LLMs for model generation tasks. Chen et al. demonstrated GPT-4's ability to generate UML class diagrams from textual descriptions~\cite{gpt4_uml}, while others have investigated LLM-based requirements-to-models transformations~\cite{llm_requirements}. A critical limitation of all these approaches is that LLMs produce \emph{textual representations} of models (e.g., PlantUML, Mermaid, or XMI fragments) that must be manually imported into EA tools. The LLM is confined to the text generation action space; it cannot instantiate model elements, create packages, or establish connectors in a professional EA tool repository. None of these works enables LLMs to directly manipulate professional EA tool repositories, which is the capability gap that MCP addresses.

\subsection{Model Context Protocol (MCP)}

MCP is an open protocol standardized by Anthropic that defines how applications can provide context and tools to LLMs in a vendor-agnostic manner~\cite{mcp_spec}. MCP follows a client-server architecture where MCP Servers expose three primitives: \emph{resources} (readable data), \emph{tools} (executable functions with typed parameters and structured responses), and \emph{prompts} (reusable templates) to MCP Clients (LLM hosts). The protocol enforces JSON Schema-based parameter validation and structured response formats, ensuring that tool invocations are deterministic, auditable, and type-safe.

In the architecture modeling domain, MCP presents a natural fit: EA tools maintain complex model repositories with rich APIs for creating packages, elements, diagrams, and connectors. By wrapping these APIs as MCP tools, LLMs gain the ability to programmatically construct architecture models through structured, auditable tool calls. To the best of our knowledge, this work represents the first application of MCP to the EA tooling domain.

\subsection{DoDAF Automation and Its Limitations}

Automating DoDAF model construction has been a long-standing goal in defense systems engineering. MDA-based approaches~\cite{mda_dodaf} transform Platform-Independent Models (PIMs) to DoDAF views through predefined mapping rules, but they require formally structured PIMs as input---a prerequisite rarely met in early-stage requirements when concepts are still being explored and formalized. Ontology-based reasoning~\cite{ontology_dodaf} can automate viewpoint selection through semantic inference, but the ontologies themselves must be manually curated and cannot accommodate novel requirement patterns that fall outside their predefined knowledge bases. Template-based generation~\cite{template_dodaf} produces views from structured databases but offers no flexibility when requirements deviate from pre-designed templates.

The common failure mode of these approaches is their \emph{brittleness to input variability}: they work well when inputs match their expected patterns but break silently or produce incorrect results when faced with the ambiguous, incomplete, and evolving requirements typical of early concept development. The integration of LLMs with DoDAF modeling introduces a paradigm shift: instead of requiring formal specifications, the system accepts natural language as input and performs semantic interpretation that bridges the gap between informal requirements and formal architectural constructs. When combined with MCP, this semantic capability is extended with direct tool-manipulation capability, closing both the semantic gap and the action space gap simultaneously.

\section{Methodology}
\label{sec:methodology}

\subsection{Overall Framework Architecture}

Our proposed framework, illustrated in Fig.~\ref{fig:architecture}, employs a multi-agent workflow comprising four specialized agents, each responsible for a distinct phase of the generation process:

\begin{enumerate}
    \item \textbf{Requirement Agent}: Parses the natural language requirement document and performs architectural decomposition---identifying operational nodes, system components, information flows, activities, and constraints. The output is a structured requirement analysis report that serves as input to the Planner Agent.
    \item \textbf{Planner Agent}: Based on the requirement analysis, determines which DoDAF viewpoints to generate, plans the hierarchical package structure, and selects appropriate SysML diagram types. Implements the formalized viewpoint selection mechanism described in Section~\ref{sec:viewpoint_selection}.
    \item \textbf{Executer Agent (MCP Client)}: Executes the generation plan by invoking MCP tools on the EAMCP Server to create packages, elements, diagrams, connectors, and layouts in the EA model repository. This agent bridges the action space between LLM reasoning and EA tool operations.
    \item \textbf{Verifier Agent}: A secondary LLM-based reviewer that validates generated models against completeness and consistency criteria, identifies errors, and triggers corrective actions by requesting the Executer Agent to revise specific elements. This agent is critical for ensuring semantic correctness beyond mere tool-call reliability.
\end{enumerate}

These agents operate sequentially with a feedback loop: the Verifier can request the Executer to revise specific model elements when quality issues are detected, forming a closed-loop quality assurance mechanism. The complete data flow is: Requirement Document $\rightarrow$ Requirement Agent $\rightarrow$ Structured Analysis $\rightarrow$ Planner Agent $\rightarrow$ Viewpoint \& Diagram Plan $\rightarrow$ Executer Agent (MCP Client) $\rightarrow$ MCP Tool Calls $\rightarrow$ EAMCP Server $\rightarrow$ EA Model Repository $\rightarrow$ Verifier Agent $\rightarrow$ Quality Report $\rightarrow$ (if issues found) Executer Agent for corrections.

\begin{figure}[t]
\centering
\framebox[\columnwidth]{\parbox{0.9\columnwidth}{\centering\vspace{1.5em}
\textbf{Multi-Agent Workflow for Automated DoDAF Generation}\\[0.5em]
\footnotesize
Requirement Document\\[0.2em]
$\downarrow$\\[0.2em]
\textbf{Requirement Agent} --- NLP parsing, architectural decomposition\\[0.2em]
$\downarrow$\\[0.2em]
\textbf{Planner Agent} --- Viewpoint selection (Algorithm~1), package planning\\[0.2em]
$\downarrow$\\[0.2em]
\textbf{Executer Agent} (MCP Client) --- MCP tool calls to EAMCP Server\\[0.2em]
$\downarrow$\\[0.2em]
\textbf{EAMCP Server} --- EA Model Repository operations\\[0.2em]
$\downarrow$\\[0.2em]
\textbf{Verifier Agent} --- Completeness \& consistency check\\[0.2em]
$\downarrow$ (if issues found)\\[0.2em]
\textbf{Executer Agent} --- Corrective revisions\\[0.2em]
$\downarrow$\\[0.2em]
Final Architecture Package
\vspace{1.5em}}}
\caption{Multi-agent workflow for automated DoDAF architecture generation. The Verifier Agent provides feedback to the Executer for iterative refinement, forming a closed-loop quality assurance mechanism.}
\label{fig:architecture}
\end{figure}

\subsection{MCP Tool Design and Semantic Illustration}

The EAMCP Server exposes 11 MCP tools, organized by function category as shown in Table~\ref{tab:mcp_tools}. Each tool accepts typed JSON parameters following a predefined JSON Schema and returns structured responses, ensuring deterministic and auditable tool invocations.

\begin{table}[h]
\centering
\caption{MCP Tools Exposed by EAMCP Server}
\label{tab:mcp_tools}
\begin{tabular}{p{3.5cm}p{4.5cm}}
\toprule
\textbf{Tool Name} & \textbf{Function} \\
\midrule
get\_root\_packages & Retrieve root-level model packages \\
get\_current\_package & Get active package context \\
create\_or\_update\_package & Create/modify package structure \\
create\_or\_update\_elements & Create model elements in EA \\
create\_or\_update\_diagram & Generate SysML/DoDAF diagrams \\
create\_or\_update\_connectors & Create relationships between elements \\
place\_elements\_on\_diagram & Position elements on diagram canvas \\
layout\_connectors & Auto-route connector paths \\
open\_diagrams & Activate diagrams for viewing \\
reload\_diagrams & Refresh diagram rendering \\
get\_diagram\_image & Export diagram as image \\
\bottomrule
\end{tabular}
\end{table}

To illustrate the semantic translation capability enabled by MCP, consider a concrete example. Given the natural language requirement ``The Emergency Command Center coordinates all responding agencies,'' traditional approaches would require a human architect to manually create a SysML Block element in EA. In our framework, the Executer Agent invokes the \texttt{create\_or\_update\_elements} tool with the following MCP request:

{\footnotesize
\begin{verbatim}
{
  "method": "tools/call",
  "params": {
    "name": "create_or_update_elements",
    "arguments": {
      "packageId": "pkg_smartcity_sv1",
      "elements": [{
        "name": "Emergency Command Center",
        "type": "Class",
        "stereotype": "SysML::Block",
        "notes": "Coordinates all responding agencies
                   including Police, Fire, Medical,
                   and Municipal Services units"
      }]
    }
  }
}
\end{verbatim}}

The EAMCP Server processes this request by directly creating a SysML Block element in the EA repository's underlying database and returns a structured response:

{\footnotesize
\begin{verbatim}
{
  "status": "success",
  "elementId": "el_047",
  "name": "Emergency Command Center",
  "type": "Class",
  "stereotype": "SysML::Block",
  "packagePath": "/Model/SmartCity/SV/SV-1",
  "createdAt": "2025-06-10T14:32:17Z"
}
\end{verbatim}}

This round-trip demonstrates the key capability enabled by MCP: the LLM does not generate PlantUML code for the user to manually import; instead, it directly instantiates architectural constructs in the EA model repository through structured, type-safe tool calls. The entire 184-tool-call pipeline across three rounds maintained this pattern without a single failure.

\subsection{Viewpoint Selection Mechanism}
\label{sec:viewpoint_selection}

The LLM agent does not attempt to generate all 52 DoDAF model products. Instead, it employs a context-aware viewpoint selection mechanism based on semantic keyword-weight mapping. Given a requirement document $R$, the Planner Agent computes a relevance score $S(v_i)$ for each DoDAF viewpoint $v_i$ as follows:

\begin{equation}
\label{eq:viewpoint}
S(v_i) = \sum_{k \in K(R)} w(v_i, k) \cdot f(k, R)
\end{equation}

\noindent where $K(R)$ is the set of domain keywords extracted from requirement $R$, $w(v_i, k)$ is the predefined semantic weight between viewpoint $v_i$ and keyword $k$, and $f(k, R)$ is the frequency-importance score of keyword $k$ in $R$. A viewpoint $v_i$ is selected if $S(v_i) \geq \theta$, where $\theta$ is a configurable threshold (empirically set to 0.15 in our experiments).

The semantic weight mapping $w(v_i, k)$ is derived from the following correspondence between SysML diagram types and DoDAF views:

\begin{itemize}
    \item \textbf{Block Definition Diagram (BDD)} $\rightarrow$ OV-1 (High-Level Operational Concept) and SV-1 (Systems Interface Description)
    \item \textbf{Internal Block Diagram (IBD)} $\rightarrow$ SV-4 (Systems Functionality) and SV-5 (Operational Activity to Systems Function Traceability)
    \item \textbf{Use Case Diagram (UC)} $\rightarrow$ OV-5 (Operational Activity Model)
    \item \textbf{Activity Diagram (ACT)} $\rightarrow$ OV-6c (Event-Trace Description)
    \item \textbf{Sequence Diagram (SD)} $\rightarrow$ OV-6c and SV-6 (Systems Resource Flow)
    \item \textbf{Requirement Diagram (REQ)} $\rightarrow$ SV-7 (Systems Measures Matrix)
    \item \textbf{State Machine Diagram (STM)} $\rightarrow$ OV-6a (Operational Rules) and SV-10a (Systems Rules)
\end{itemize}

The viewpoint selection procedure is formalized in Algorithm~\ref{alg:viewpoint}. Unlike prior approaches that hard-code viewpoint selection rules, this mechanism is configurable via the weight matrix $W$ and threshold $\theta$, enabling adaptation to different domains and architecture frameworks.

\begin{algorithm}[t]
\caption{Context-Aware Viewpoint Selection}
\label{alg:viewpoint}
\begin{algorithmic}[1]
\REQUIRE Requirement document $R$, weight matrix $W$, threshold $\theta$
\ENSURE Selected viewpoint set $V_{sel}$
\STATE $K \leftarrow \text{ExtractKeywords}(R)$ \COMMENT{extract domain keywords}
\STATE $V_{sel} \leftarrow \emptyset$
\FOR{each viewpoint $v_i$ in DoDAF viewpoint set $\mathcal{V}$}
    \STATE $S(v_i) \leftarrow 0$
    \FOR{each keyword $k_j$ in $K$}
        \STATE $S(v_i) \leftarrow S(v_i) + w(v_i, k_j) \cdot f(k_j, R)$
    \ENDFOR
    \IF{$S(v_i) \geq \theta$}
        \STATE $V_{sel} \leftarrow V_{sel} \cup \{v_i\}$
    \ENDIF
\ENDFOR
\RETURN $V_{sel}$
\end{algorithmic}
\end{algorithm}

\subsection{End-to-End Generation Pipeline}

The generation proceeds through six phases, each executed by the corresponding agent:

\textbf{Phase 1: Requirement Analysis (Requirement Agent).} The LLM reads the natural language requirement document and performs architectural decomposition: identifying operational nodes, system components, information flows, activities, and constraints. The output is a structured analysis report.

\textbf{Phase 2: Viewpoint Selection and Planning (Planner Agent).} The Planner Agent applies Algorithm~\ref{alg:viewpoint} to determine which DoDAF viewpoints to generate and plans the hierarchical package structure with sub-packages for DoDAF operational views and SysML structural/behavioral views.

\textbf{Phase 3: Package Structure Creation (Executer Agent).} The Executer creates the package hierarchy in the EA repository through MCP calls to \texttt{create\_or\_update\_package}.

\textbf{Phase 4: Element and Diagram Generation (Executer Agent).} For each selected viewpoint, the Executer creates model elements (actors, components, nodes, interfaces) via \texttt{create\_or\_update\_elements}, generates diagrams via \texttt{create\_or\_update\_diagram}, and establishes connectors via \texttt{create\_or\_update\_connectors}.

\textbf{Phase 5: Layout and Refinement (Executer Agent).} Elements are positioned via \texttt{place\_elements\_on\_diagram}, and connectors are auto-routed via \texttt{layout\_connectors} for readability.

\textbf{Phase 6: Quality Verification (Verifier Agent).} The Verifier Agent examines the generated model for completeness and consistency. If issues are detected, it triggers corrective actions by requesting the Executer to revise specific elements, forming a closed-loop quality assurance loop.

\subsection{Quality Verification}
\label{sec:quality_verification}

The Verifier Agent evaluates generated models along four dimensions, each producing quantifiable metrics:

\begin{itemize}
    \item \textbf{Completeness}: Are all key DoDAF viewpoints addressed? Are essential system components represented? Measured as the ratio of covered viewpoints to recommended viewpoints.
    \item \textbf{Consistency}: Do elements referenced across multiple diagrams maintain consistent names, types, and relationships? Measured as the number of cross-diagram inconsistency issues detected and corrected.
    \item \textbf{Reliability}: What proportion of MCP tool calls succeed? Measured as the tool call success rate.
    \item \textbf{Semantic Correctness}: Does the generated architecture correctly reflect the original requirements? Measured through domain expert review on a graded rubric (0--100\%).
\end{itemize}

It is critical to distinguish between these dimensions: reliability (tool calls succeed) and semantic correctness (generated content accurately reflects requirements) are orthogonal. A pipeline can achieve 100\% reliability while producing semantically incorrect models if, for example, it creates elements with wrong stereotypes or omits key relationships. The Verifier Agent specifically targets the correctness dimension by checking cross-viewpoint consistency, naming conventions, and requirement-to-element traceability, closing the gap between ``the pipe works'' and ``the content is right.''

\section{Case Study: Smart City Emergency Command System}
\label{sec:casestudy}

\subsection{System Description}

The Smart City Emergency Command System is a civil emergency management platform that coordinates multi-agency response to urban incidents. The system requirements encompass: \textbf{Situation Awareness} (real-time monitoring of urban sensors, surveillance cameras, and public reporting channels); \textbf{Incident Triage} (automated classification and prioritization of incoming emergency reports); \textbf{Resource Dispatch} (coordination of police, fire, medical, and municipal response units); \textbf{Command \& Control} (multi-level command hierarchy with information sharing and decision support); \textbf{Communication} (interoperable communication across agencies with different equipment and protocols); and \textbf{Post-Incident Analysis} (data collection for after-action review and process improvement).

This system was chosen as a validation case because it represents a complex system-of-systems with rich operational and system-level interactions, making it suitable for demonstrating the framework's capability to generate multi-viewpoint DoDAF architectures.

\subsection{Generated Architecture Models}

The LLM agent, given the natural language requirements, autonomously selected and generated architecture artifacts across three experimental rounds. Table~\ref{tab:viewpoints} summarizes the viewpoints generated.

\begin{table}[h]
\centering
\caption{DoDAF Viewpoints Generated Across Three Rounds}
\label{tab:viewpoints}
\begin{tabular}{lccc}
\toprule
\textbf{Viewpoint} & \textbf{Round~1} & \textbf{Round~2} & \textbf{Round~3} \\
\midrule
AV-1 (Overview \& Summary) & \checkmark & \checkmark & -- \\
AV-2 (Integrated Dictionary) & \checkmark & -- & -- \\
OV-1 (High-Level Concept) & \checkmark & \checkmark & \checkmark \\
OV-2 (Resource Flow) & \checkmark & \checkmark & \checkmark \\
OV-4 (Organizational) & \checkmark & -- & \checkmark \\
OV-5 (Operational Activity) & \checkmark & \checkmark & \checkmark \\
OV-6a (Operational Rules) & \checkmark & -- & -- \\
OV-6b (State Transition) & \checkmark & -- & -- \\
OV-6c (Event-Trace) & \checkmark & -- & -- \\
SV-1 (Systems Interface) & \checkmark & \checkmark & \checkmark \\
SV-2 (Resource Flow) & \checkmark & -- & -- \\
SV-3 (Systems Matrix) & \checkmark & -- & -- \\
SV-4 (Systems Functionality) & \checkmark & \checkmark & -- \\
SV-5 (Activity Traceability) & \checkmark & -- & -- \\
SV-6 (Data Exchange) & \checkmark & -- & -- \\
SV-7 (Measures) & \checkmark & -- & -- \\
SV-8 (Evolution) & \checkmark & -- & -- \\
SV-9 (Technology Forecast) & \checkmark & -- & -- \\
SV-10a (Systems Rules) & \checkmark & -- & -- \\
SV-10b (State Transition) & \checkmark & -- & -- \\
SV-10c (Event-Trace) & \checkmark & -- & -- \\
\bottomrule
\end{tabular}
\end{table}

Round~1 exhibited the most comprehensive viewpoint coverage, generating 21 distinct DoDAF views including all three SV-10 variants. Rounds~2 and~3 demonstrated more focused selections, converging on the core operational and systems views (OV-1, OV-2, OV-5, SV-1, SV-4) that are most critical for initial architecture understanding.

\subsection{Illustrative Generated Diagrams}

Fig.~\ref{fig:ov1} and Fig.~\ref{fig:sv1} illustrate representative SysML diagrams generated by the framework and exported from the EA model repository. These are not manually drawn but are the actual output of the automated pipeline.

\begin{figure}[t]
\centering
\framebox[\columnwidth]{\parbox{0.9\columnwidth}{\centering\vspace{1.2em}
\textbf{OV-1: High-Level Operational Concept}\\[0.4em]
\footnotesize
Emergency Command Center $\leftrightarrow$ Police Unit (Incident Assignment)\\
Emergency Command Center $\leftrightarrow$ Fire Rescue Unit (Resource Request)\\
Emergency Command Center $\leftrightarrow$ Medical Response Unit (Casualty Report)\\
Emergency Command Center $\leftrightarrow$ Municipal Services (Infrastructure Status)\\
Sensor Network $\rightarrow$ Emergency Command Center (Real-time Data Feed)\\
Public Reporting Channel $\rightarrow$ Incident Triage System (Emergency Reports)\\
\vspace{1.2em}}}
\caption{OV-1 High-Level Operational Concept diagram generated by the framework, showing operational nodes and their information exchange relationships.}
\label{fig:ov1}
\end{figure}

\begin{figure}[t]
\centering
\framebox[\columnwidth]{\parbox{0.9\columnwidth}{\centering\vspace{1.2em}
\textbf{SV-1: Systems Interface Description}\\[0.4em]
\footnotesize
[Command \& Control System] ---(REST/SOAP)--- [Resource Dispatch Module]\\
[Sensor Integration Platform] ---(MQTT)--- [Situation Awareness Dashboard]\\
[Communication Gateway] ---(SIP/RTP)--- [Interagency Comm Module]\\
[Incident Triage Engine] ---(Kafka)--- [Command \& Control System]\\
[GIS Mapping Service] ---(WMS)--- [Situation Awareness Dashboard]\\
\vspace{1.2em}}}
\caption{SV-1 Systems Interface Description diagram generated by the framework, showing system components and their interface protocols.}
\label{fig:sv1}
\end{figure}

\section{Evaluation}
\label{sec:evaluation}

We evaluate the framework along three distinct dimensions---\emph{efficiency} (how fast), \emph{reliability} (whether tools work), and \emph{semantic correctness} (whether the content is right)---to avoid conflating pipeline throughput with output quality.

\subsection{Experimental Setup}

All experiments were conducted using DeepSeek-V4 as the LLM backend, accessed through a unified API interface. To ensure reproducibility, all generation tasks used temperature~$=$~0.3, while the weight matrix generation (Algorithm~\ref{alg:viewpoint}) used temperature~$=$~0 for fully deterministic output. The MCP Server (EAMCP) was deployed locally, connecting to Sparx Systems Enterprise Architect. Each experimental round started with an empty EA project and the same natural language requirement document describing the Smart City Emergency Command System. Two domain experts (senior system architects with 10+ years of DoDAF experience at the authors' institution) independently assessed the semantic correctness of post-verification models.

\subsection{Efficiency Analysis}

Table~\ref{tab:efficiency} presents the efficiency metrics across three rounds.

\begin{table}[h]
\centering
\caption{Generation Efficiency Across Experimental Rounds}
\label{tab:efficiency}
\begin{tabular}{lcccc}
\toprule
\textbf{Metric} & \textbf{Round~1} & \textbf{Round~2} & \textbf{Round~3} & \textbf{Avg.} \\
\midrule
Total Duration & 4~min~20~s & 5~min~35~s & 4~min~35~s & 4~min~50~s \\
MCP Tool Calls & 66 & 70 & 48 & 61.3 \\
Diagrams Generated & 11 & 10 & 6 & 9.0 \\
Package Operations & 9 & 8 & 3 & 6.7 \\
Element Operations & 11 & 17 & 9 & 12.3 \\
Connector Operations & 9 & 12 & 7 & 9.3 \\
Layout Operations & 11 & 10 & 7 & 9.3 \\
\bottomrule
\end{tabular}
\end{table}

The average end-to-end automated pipeline time of 4~min~50~s---which includes LLM-driven generation, Verifier Agent detection, and automated correction of detected issues---represents a dramatic improvement over manual modeling. Based on the authors' professional experience in defense system architecture development, an experienced architecture engineer would require approximately 4--8~hours to manually produce a comparable set of DoDAF views in Enterprise Architect, including package structure creation, element definition, diagram drawing, connector establishment, and layout refinement. This represents a speedup of approximately 50--100$\times$ for the automated pipeline (generation $+$ detection $+$ auto-fix), noting that the human reviewer is repositioned to a supervisory role rather than being eliminated from the workflow.

Round~2 exhibited the highest element creation count (17 calls) and longest duration (5~min~35~s), correlating with more detailed model construction. Round~3 achieved the fastest generation (4~min~35~s) with the fewest tool calls (48), reflecting more efficient viewpoint selection after the Planner Agent refined its strategy based on the requirement analysis.

\subsection{Reliability Analysis}

A key finding is the zero-failure reliability: across all three rounds and 184 total MCP tool calls, not a single call failed (100\% success rate). This demonstrates that the MCP tool chain is robust and that the LLM generates syntactically correct tool parameters. Several factors contribute to this reliability: (1) deterministic tool interfaces with well-typed JSON parameters and structured responses reduce ambiguity; (2) contextual grounding---the LLM retrieves current EA project state before making modifications, ensuring operations target valid contexts; and (3) sequential dependency management---the LLM respects the ordering of packages before elements, elements before diagrams, and diagrams before connectors.

\subsection{Semantic Correctness Analysis}

We emphasize that \emph{reliability $\neq$ correctness}: 100\% tool call success only proves that ``the pipeline works,'' not that ``the content is right.'' To evaluate semantic correctness, we employ both the Verifier Agent (automated) and domain expert review (manual). Table~\ref{tab:quality} presents the quality assessment results.

\begin{table}[h]
\centering
\caption{Quality Assessment Results}
\label{tab:quality}
\begin{tabular}{lcccc}
\toprule
\textbf{Metric} & \textbf{Round~1} & \textbf{Round~2} & \textbf{Round~3} & \textbf{Avg.} \\
\midrule
Tool Call Success Rate & 100\% & 100\% & 100\% & 100\% \\
Pre-Verification Issues Detected & 5 & 4 & 2 & 3.7 \\
Post-Verification Issues Remaining & 1 & 1 & 0 & 0.7 \\
Issues Corrected by Verifier & 4 & 3 & 2 & 3.0 \\
Expert Correctness Score & 88\% & 92\% & 97\% & 92.3\% \\
\bottomrule
\end{tabular}
\end{table}

The Verifier Agent detected an average of 3.7 consistency issues per round before verification. These issues included: naming inconsistencies across diagrams (e.g., ``Emergency\_Cmd\_Center'' in SV-1 vs.\ ``Emergency Command Center'' in OV-1), missing connectors between related elements (e.g., Police Unit not connected to Command Center in OV-2 despite appearing in OV-1), and incomplete viewpoint coverage where a recommended view was not generated. After Verifier-triggered corrections, the average number of remaining issues dropped to 0.7 per round.

Two domain experts independently assessed the semantic correctness of the post-verification models, scoring an average of 92.3\% across three rounds. The most common remaining issues were suboptimal element naming conventions and minor omissions in non-core viewpoints (e.g., SV-9 Technology Forecast generated with incomplete technology timelines). The improvement from Round~1 (88\%) to Round~3 (97\%) reflects the iterative refinement of the Planner Agent's viewpoint selection and the Verifier Agent's rule set across experimental rounds.

\subsection{Ablation Study: Impact of the Verifier Agent}
\label{sec:ablation}

To quantify the contribution of the Verifier Agent, we compare the quality of models generated with and without the verification step. Table~\ref{tab:ablation} presents the ablation results.

\begin{table}[h]
\centering
\caption{Ablation Study: With vs.\ Without Verifier Agent}
\label{tab:ablation}
\begin{tabular}{lcc}
\toprule
\textbf{Metric} & \textbf{Without Verifier} & \textbf{With Verifier} \\
\midrule
Avg. Consistency Issues & 3.7 & 0.7 \\
Issue Reduction Rate & -- & 81\% \\
Expert Correctness Score & 78.3\% & 92.3\% \\
\bottomrule
\end{tabular}
\end{table}

Without the Verifier, the average number of cross-diagram consistency issues per round was 3.7; with the Verifier, this dropped to 0.7---an 81\% reduction. Expert correctness scores without verification averaged 78.3\% (vs.\ 92.3\% with verification), confirming that the Verifier Agent provides substantial quality improvement beyond mere tool-call reliability. This ablation demonstrates that the closed-loop verification mechanism is not an optional addition but a necessary component for producing semantically correct architecture models.

\subsection{Tool Call Distribution Analysis}

Table~\ref{tab:tool_distribution} breaks down the tool call distribution by function category.

\begin{table}[h]
\centering
\caption{Tool Call Distribution by Function Category}
\label{tab:tool_distribution}
\begin{tabular}{lccc}
\toprule
\textbf{Category} & \textbf{Round~1} & \textbf{Round~2} & \textbf{Round~3} \\
\midrule
Context Operations & 2 & 2 & 2 \\
Package Management & 9 & 8 & 3 \\
Element Creation & 11 & 17 & 9 \\
Diagram Creation & 11 & 10 & 6 \\
Connector Creation & 9 & 12 & 7 \\
Element Placement & 11 & 10 & 7 \\
Layout \& Rendering & 12 & 10 & 8 \\
Diagram Export & 0 & 0 & 6 \\
\midrule
\textbf{Total} & \textbf{66} & \textbf{70} & \textbf{48} \\
\bottomrule
\end{tabular}
\end{table}

Element creation and diagram creation consistently account for the largest share of tool calls (approximately 33\% combined), followed by layout operations (17\%) and connector management (16\%). Context operations are consistently minimal at 2~calls per round, serving only as initialization steps.

\subsection{Viewpoint Coverage Analysis}

Viewpoint coverage varied across rounds, reflecting the LLM's autonomous selection strategy. Round~1 generated the most comprehensive coverage (21 distinct DoDAF views), while Rounds~2 and~3 converged on a core set of 5--6 views. This variability highlights both a strength and a limitation of the autonomous approach: the LLM can adapt its coverage based on requirement interpretation, but the selection is not guaranteed to be consistent across runs. For production use, a hybrid approach---LLM-driven generation with human-specified mandatory viewpoints---may be preferable. The formalized viewpoint selection mechanism (Algorithm~\ref{alg:viewpoint}) provides a foundation for such hybrid approaches by making the selection logic transparent and configurable.

\subsection{Discussion}

The experimental results validate the feasibility of LLM-driven automated DoDAF architecture generation, with important nuances.

\textbf{Efficiency vs.\ Quality Trade-off:} The 4~min~50~s end-to-end automated pipeline time (generation + Verifier detection + auto-correction) is impressive, but speed alone is insufficient. The Verifier Agent's contribution (81\% reduction in consistency issues, 92.3\% expert correctness score) demonstrates that closed-loop quality assurance is essential for producing usable architecture models. The automated detection-and-fix cycle is a critical component of the pipeline: without it, the human reviewer would need to manually identify and resolve each inconsistency, negating the time savings.

\textbf{Reliability $\neq$ Correctness:} Our results clearly separate these two dimensions, which prior work has often conflated. The 100\% tool call success rate proves pipeline reliability; the 92.3\% expert correctness score proves semantic quality. Both metrics are necessary but individually insufficient for evaluating an architecture generation framework.

\textbf{Limitations:} (1) The quality of generated models depends on LLM capability and may vary across different LLM backends. (2) Auto-layout diagrams may require manual aesthetic refinement before formal review. (3) The Smart City case study, while representative, does not cover classified military scenarios where requirement sensitivity constrains LLM usage. (4) The expert correctness assessment involved two reviewers; a larger-scale user study would strengthen the evaluation. (5) The comparison with manual modeling is based on estimated rather than measured times; a controlled user study is planned as future work.

\section{Conclusion}
\label{sec:conclusion}

This paper presented an LLM-driven MCP framework for automated DoDAF-compliant architecture model generation. By integrating multi-agent LLM reasoning with MCP-enabled EA tool orchestration, the framework achieves end-to-end automation from natural language requirements to complete architecture packages. Three contributions distinguish this work from prior approaches. First, we established the first integrated LLM-MCP-DoDAF pipeline, demonstrating that MCP bridges the action space gap that has limited prior LLM-based modeling to textual code generation. Second, we formalized a context-aware viewpoint selection mechanism (Algorithm~\ref{alg:viewpoint}) that makes the LLM's architectural decision-making transparent and configurable. Third, we demonstrated the critical role of the Verifier Agent through an ablation study (81\% reduction in consistency issues) and separated the evaluation dimensions of efficiency, reliability, and semantic correctness---revealing that the latter two are orthogonal and must be assessed independently.

Experimental evaluation using a Smart City Emergency Command System case study with DeepSeek-V4 demonstrated an average end-to-end automated pipeline time of 4~min~50~s (generation + Verifier detection + auto-correction, approximately 50--100$\times$ speedup over manual modeling), 100\% tool-call reliability, and a post-verification semantic correctness score of 92.3\%.

Future work will address several directions. First, we plan to conduct controlled user studies comparing framework-generated architectures with manually produced ones, measuring both time and quality differences rigorously. Second, we aim to evaluate the framework with additional LLM backends (GPT-4, Claude) and across diverse domains beyond smart city systems. Third, we will investigate the integration of formal verification techniques to provide mathematical guarantees on model consistency beyond the current heuristic quality checks. Fourth, we plan to explore the applicability of this framework to other architecture frameworks such as MODAF, NAF, and TOGAF.

\section*{Data Availability}
The natural language requirement document for the Smart City case study, the complete experimental data (tool call traces, Verifier reports, and ablation study results across 10~rounds), and the analyzed metrics are publicly available at \url{https://github.com/liu-weigang/DoDAF-MCP-Experiment}. The EAMCP (Enterprise Architect MCP Server) source code is not publicly released due to institutional restrictions.

\begin{thebibliography}{99}

\bibitem{DoDAF2010}
DoD Deputy Chief Information Officer, ``The DoDAF Architecture Framework Version 2.02,'' U.S. Department of Defense, Aug. 2010. [Online]. Available: \url{https://dodcio.defense.gov/Library/DoD-Architecture-Framework/}

\bibitem{llm_survey}
W.~X. Zhao, K. Zhou, J. Li, \textit{et al.}, ``A Survey of Large Language Models,'' \textit{arXiv preprint}, arXiv:2303.18223, 2023.

\bibitem{mcp_spec}
Anthropic, ``Model Context Protocol Specification,'' 2024. [Online]. Available: \url{https://modelcontextprotocol.io/}

\bibitem{uml_nlp}
D. Yue, L. Briand, and Y. Labiche, ``aToucan: An Automated Framework to Derive UML Analysis Models from Use Case Models,'' \textit{ACM Trans. Softw. Eng. Methodol.}, vol.~24, no.~3, pp.~1--52, 2015.

\bibitem{deep_modeling}
S. Ahmed, A. Ahmed, and N.~H.~M. Sani, ``Automated UML Class Diagram Generation from Textual Requirements Using Deep Learning,'' in \textit{Proc. IEEE Int. Conf. Comput. Commun. Eng. (ICCCE)}, 2021, pp.~206--211.

\bibitem{gpt4_uml}
J. Chen, Y. Huang, and P. Li, ``Can GPT-4 Generate UML Diagrams from Natural Language? A Preliminary Study,'' in \textit{Proc. IEEE/ACM Int. Conf. Autom. Softw. Eng. (ASE)}, 2024, pp.~1--6.

\bibitem{llm_requirements}
M. Vierhauser, R. Rabiser, and P. Gr\"{u}nbacher, ``Requirements Engineering in the Age of Large Language Models: A Systematic Mapping Study,'' \textit{Inf. Softw. Technol.}, vol.~168, p.~107369, 2024.

\bibitem{mda_dodaf}
S. Mittal, ``Extending DoDAF to Allow Integrated DEVS-Based Modeling and Simulation,'' \textit{J. Defense Model. Simul.}, vol.~3, no.~2, pp.~95--123, 2006.

\bibitem{ontology_dodaf}
P.~C.~G. Costa, K.~B. Laskey, and T. Lukasiewicz, ``Uncertainty Representation and Reasoning in the Semantic Web,'' in \textit{Semantic Web Engineering in the Knowledge Society}. IGI Global, 2009, pp.~315--340.

\bibitem{template_dodaf}
L. Zhang, Y. Luo, and F. Tao, ``Cloud Manufacturing: A New Manufacturing Paradigm,'' \textit{Enterprise Inf. Syst.}, vol.~8, no.~2, pp.~167--187, 2014.

\bibitem{ai_se}
D.~M.~A. Hussain and M.~A. Ahmed, ``Artificial Intelligence in Systems Engineering: A Systematic Literature Review,'' \textit{IEEE Access}, vol.~12, pp.~34567--34589, 2024.

\bibitem{sysml_spec}
Object Management Group, ``OMG Systems Modeling Language (SysML) Version 1.6,'' Nov. 2019. [Online]. Available: \url{https://www.omg.org/spec/SysML/}

\bibitem{ea_tool}
Sparx Systems, ``Enterprise Architect User Guide,'' 2024. [Online]. Available: \url{https://sparxsystems.com/enterprise_architect_user_guide/}

\bibitem{deepseek}
DeepSeek-AI, ``DeepSeek-V4 Technical Report,'' 2025. [Online]. Available: \url{https://deepseek.com/}

\bibitem{smartcity}
C. Yin, Z. Xiong, H. Chen, J. Wang, D. Cooper, and B. David, ``A Literature Survey on Smart Cities,'' \textit{Sci. China Inf. Sci.}, vol.~58, no.~10, pp.~1--18, 2015.

\end{thebibliography}

\begin{IEEEbiography}{Weigang Liu}
received the B.S. and M.S. degrees from Tongji University, Shanghai, China. He is currently an engineer at the Network and Communication Research Institute of China Electronics Technology Group Corporation. His research interests include tactical data links, enterprise architecture, and AI-driven systems engineering. He is a recipient of the CETC Outstanding Youth Award.
\end{IEEEbiography}

\begin{IEEEbiography}{Xing Zhou}
is an engineer at the Network and Communication Research Institute of China Electronics Technology Group Corporation. His research interests include communication systems and network architecture.
\end{IEEEbiography}

\begin{IEEEbiography}{Zhongjie Qi}
is an engineer at the Network and Communication Research Institute of China Electronics Technology Group Corporation. His research focuses on tactical data link message processing and information system architecture.
\end{IEEEbiography}

\begin{IEEEbiography}{Zhenfeng Yan}
received the Ph.D. degree from Tongji University. He is currently a senior engineer at the Network and Communication Research Institute of China Electronics Technology Group Corporation. His research interests include C4ISR systems, tactical data links, and system-of-systems engineering.
\end{IEEEbiography}

\begin{IEEEbiography}{Xuepan Gao}
(Corresponding Author) is an engineer at the Network and Communication Research Institute of China Electronics Technology Group Corporation. His research interests include software-defined networking and military communication systems.
\end{IEEEbiography}

\begin{IEEEbiography}{Chuan Geng}
is an engineer at the Network and Communication Research Institute of China Electronics Technology Group Corporation. His research interests include information system integration and data link protocols.
\end{IEEEbiography}

\begin{IEEEbiography}{Haiyang Ren}
is an engineer at the Network and Communication Research Institute of China Electronics Technology Group Corporation. His research interests include embedded systems and real-time communication.
\end{IEEEbiography}

\EOD

\end{document}
